\section{Improved Junta Theorems}\label{sec:super_fried}
\subsection{Proof of Theorem~\ref{thm:super_fried}}
In this section, we prove Theorem~\ref{thm:super_fried}. As explained in the introduction, if one does not care about getting tight dependency
between the junta size $J$ and the parameter $A$, then one may use Corollary~\ref{corr:noise_stable} and Bourgain's noise sensitivity theorem~\cite{Bourgain,kindler2018gaussian}
and get a junta theorem with slightly worse parameters than those stated in Theorem~\ref{thm:super_fried}. The main goal of this section
is to get an optimal dependency of $J$ on $A$.

For the proof, we will need the notion of noisy influences of a function $f$, and we first generalize the notion of influences to real valued functions.
\begin{definition}
  Let $f\colon\power{n}\to\mathbb{R}$ be a function, and let $i\in [n]$ be a coordinate.
  We define the influence of $i$ on $f$ as $I_i[f] = \norm{\partial_i f}_2^2$, where
  we recall that $\partial_i f$ is the discrete derivative of $f$ along direction $i$.
\end{definition}
\begin{definition}
  The $\rho$-noisy influence of a coordinate $i$ for $f\colon\{0,1\}^n\to\mathbb{R}$ is defined to be
  $I_i^{(\rho)}[f] = I_i[T_{\rho} f]$.
\end{definition}

\begin{fact}\label{fact:sum_of_noisy}[\cite{ODonnell}]
  For all $f\colon\{0,1\}^n\to\{0,1\}$ we have that $\sum\limits_{i=1}^{n} I_i^{(\rho)}[f] \ll \frac{1}{1-\rho}$.
  \qed
\end{fact}


Let $C_1, C_2$ be two absolute constants, and we think of $C_2$ as much larger than $C_1$. Given $f\colon\{0,1\}^n\to\{0,1\}$ with
$\Expect{x}{s_f(x)^p}\leq A$, we denote $d = C_1\left(\frac{A}{\eps}\right)^{1/p}$, $\delta = 2^{-\frac{C_2}{2p-1} d}$.
Define the candidate set of junta coordinates as
\[
\mathcal{J} = \sett{i\in [n]}{I_i[T_{1-\frac{1}{3d}}[f]]\geq \delta},
\]
i.e.~$\mathcal{J}$ is the set of coordinates whose noisy influence at $f$ is somewhat large. Our task is to show that $f$ is $\eps$-close to a $\mathcal{J}$ junta,
i.e.~that
\begin{equation}\label{eq6}
\sum\limits_{S: S\not\subseteq \mathcal{J}}\widehat{f}(S)^2\leq \eps,
\end{equation}
and the rest of the proof is devoted towards this goal. We remark that the junta $\mathcal{J}$ is
similar in spirit to the junta in Bourgain’s original proof~\cite{Bourgain}, which is easier to see in~\cite[Section 4]{kindler2018gaussian} and
in particular in~\cite[Theorem 4.20]{kindler2018gaussian}. Indeed, both here and in~\cite[Section 4]{kindler2018gaussian}, the junta is
taken to be the set of coordinates with significant noisy influence. There are many other similarities between the argument
presented therein and the argument here. First, in both cases one considers a dyadic sort of partitioning of the Fourier coefficients:
in~\cite[Section 4]{kindler2018gaussian} it is achieved via considering the difference in the noise sensitivity of $f$ at different
noise rates (see~\cite[Theorem 4.16]{kindler2018gaussian}). In the current proof, we perform explicit dyadic partitioning and also
accompany it by applications of suitable noise operators (the operators $S_k$, $S_k'$ and $S_k''$ below). Second, both proofs also use some relation between the noise sensitivity / moments
of the sensitivity and the level $1$ weight of the function $f$ (and its restrictions). While the technical execution of the proofs is
different, it is hard to pinpoint the high level difference between them.
For example, the ``harsh'' degree truncations
in~\cite{kindler2018gaussian} have to be replaced by the ``softer'' type of degree truncations that are achieved by the operators $S_k$, $S_k'$ and
$S_k''$ (see also the parameter $\tilde{W}_{k,d}$) to make the argument go through. Additionally, towards the end of the argument we are able to bound the sum of noisy influences in a non-trivial way, which ultimately allows us to choose $\delta = 2^{O_p(d)}$ as opposed to
$\delta = 2^{O_{p}(d\log d)}$ (which would have resulted in slightly worse bounds as in~\cite[Theorem 4.20]{kindler2018gaussian}); see Claim~\ref{claim:bound_low_deg_inf}.


Towards proving~\eqref{eq6}, we partition the contribution of different scales of
$\card{S\cap \mathcal{J}}$ to the left hand side, and for each $k=1,\ldots,\log d$ define
\[
W_{k,d}[f] = \sum\limits_{\substack{S: \card{S}\leq d,\\ 2^{k-1}\leq \card{S\cap \bar{\mathcal{J}}}\leq 2^k}}\widehat{f}(S)^2.
\]
Thus, we may write
\begin{equation}\label{eq:super_frid1}
  \sum\limits_{S: S\not\subseteq \mathcal{J}}\widehat{f}(S)^2
  \leq
  W_{>d}[f]
  +
  \sum\limits_{k=1}^{\log d} W_{k,d}[f].
\end{equation}
The rest of the proof is devoted to bounding $W_{>d}[f]$ and $W_{k,d}[f]$ for each $k$.
From Lemma~\ref{lemma:tal_lvld_above_alpha} it follows that
\begin{equation}\label{eq:super_frid}
  W_{>d}[f]\leq O\left(\frac{A}{d^p}\right) \leq \frac{\eps}{16}
\end{equation}
for sufficiently large $C_1$. Next, we bound $W_{k,d}[f]$. Fix $k$, and consider the operator $S_k = T_{1-1/d}^{\mathcal{J}}T_{1-1/2^{k}}^{\bar{\mathcal{J}}}$,
i.e. the noise operator that applies $1/2^k$ noise on coordinates outside $\mathcal{J}$, and $1/d$ noise on coordinates of $\mathcal{J}$. Note that
\[
W_{k,d}[f] \ll W_{k,d}[S_k f],
\]
so it suffices to bound $W_{k,d}[S_k f]$. We partition $W_{k,d}[S_k f]$ according to contribution of each $j\in\bar{\mathcal{J}}$, getting:
\begin{align*}
W_{k,d}[S_k f]
=\sum\limits_{\substack{S: \card{S}\leq d,\\ 2^{k-1}\leq \card{S\cap \bar{\mathcal{J}}}\leq 2^k}} \widehat{S_k f}(S)^2
\leq \frac{1}{2^{k-1}} \sum\limits_{j\not\in\mathcal{J}}\sum\limits_{\substack{S\ni j: \card{S}\leq d,\\ 2^{k-1}\leq \card{S\cap \bar{\mathcal{J}}}\leq 2^k}} \widehat{S_k f}(S)^2
&\ll \frac{1}{2^{k-1}} \sum\limits_{j\not\in\mathcal{J}}\norm{S_k \partial_j f}_2^2\\
&\defeq \tilde{W}_{k,d}[f].
\end{align*}
Thus, we have that $W_{k,d}[f]\ll \tilde{W}_{k,d}[f]$, and the rest of the proof is devoted to upper bounding the right hand side.
Define the operators $S_{k}'$ and $S_k''$ as
\[
S_k' = T_{1-\frac{1}{3d}}^{\mathcal{J}}T_{1-\frac{1}{2^{k}}+\frac{2}{3d}}^{\bar{\mathcal{J}}},
\qquad
S_k'' = T_{1-\frac{2}{3d}}^{\mathcal{J}}T_{1-\frac{1}{2^{k}}+\frac{1}{3d}}^{\bar{\mathcal{J}}}.
\]

We establish the following claim using random restrictions and hypercontractivity:
\begin{claim}\label{claim:super_fried_main}
  For all $\tau\in (0,1)$, we have that
  \[
  2^k \tilde{W}_{k,d}[f]\leq
  C\cdot d^{1-p}\tau^{2p-1} \Expect{x}{s_f(x)^p} + C\cdot\tau^{-\frac{1}{100 d}}\sum\limits_{j\in\bar{\mathcal{J}}} \norm{S_k' \partial_j f}_2^{2+\frac{1}{100d}},
  \]
  where $C$ is an absolute constant.
\end{claim}
\begin{proof}
  Deferred to Section~\ref{sec:missing_proofs}.
\end{proof}

Given Claim~\ref{claim:super_fried_main}, one may pull out $\norm{S_k' \partial_j f}_2^{\frac{1}{100 d}}\leq \delta^{\frac{1}{200d}}$ outside
the sum, and bound the sum on the rest of the noisy influences by $O(d)$ to get a bound of the form
$2^k \tilde{W}_{k,d}[f]\ll d^{1-p}\tau^{2p-1} \Expect{x}{s_f(x)^p} + d\tau^{-\frac{1}{100 d}}\delta^{\frac{1}{200d}}$.
This idea is good enough to establish a junta theorem, yet it gives worse parameters.\footnote{One has to take
$\delta = 2^{-O(d\log d)}$, as opposed to $\delta = 2^{-O(d)}$ as we have taken. We remark that if one only cares
about getting a larger junta of size $2^{O(d\log d)}$, the proof greatly simplifies.} Instead, to get the tight result with respect
to the junta size, we will need a more careful bound on the sum of noisy influences outside $\mathcal{J}$. Specifically,
we use an improved bound of the sum of noisy influences that relates them to $W_{\ell,d}[f]$.
To state this bound we define
\[
\eps_{\ell,d} = \sum\limits_{\substack{\card{S}>d\\ 2^{\ell-1}\leq \card{S\cap \bar{\mathcal{J}}}\leq 2^{\ell}}} \widehat{f}(S)^2.
\]
 \begin{claim}\label{claim:bound_low_deg_inf}
    $\frac{1}{2^k}\sum\limits_{j\in\bar{\mathcal{J}}} \norm{S_k' \partial_j f}_2^{2}\ll \sum\limits_{\ell=1}^{\log n} 2^{-\card{\ell-k}}(W_{\ell,d}[f]+\eps_{\ell, d})$.
  \end{claim}
  \begin{proof}
    By definition,
    \[
    \frac{1}{2^k}\sum\limits_{j\in\bar{\mathcal{J}}} \norm{S_k' \partial_j f}_2^{2}
    =4\sum\limits_{S} \frac{\card{S\cap \bar{\mathcal{J}}}}{2^k} \left(1-\frac{1}{2^k}+\frac{2}{3d}\right)^{\card{S\cap \bar{\mathcal{J}}}-1}\left(1-\frac{1}{3d}\right)^{\card{S\cap \mathcal{J}}}\widehat{f}(S)^2.
    \]
    Partitioning the last sum according to $\card{S\cap \bar{\mathcal{J}}}$ and dyadically partitioning so that $2^{\ell-1}\leq \card{S\cap \bar{\mathcal{J}}} < 2^{\ell}$, we have that the last
    expression is at most
    \[
    \ll
    \sum\limits_{\ell=1}^{\log n} 2^{\ell-k} \left(1-\frac{1}{2^k}+\frac{2}{3d}\right)^{2^{\ell-1}}(W_{\ell,d}[f]+\eps_{\ell, d})
    \ll
    \sum\limits_{\ell=1}^{\log n} 2^{-\card{\ell-k}}(W_{\ell,d}[f]+\eps_{\ell, d}).\qedhere
    \]
  \end{proof}

    Combining Claims~\ref{claim:super_fried_main} and~\ref{claim:bound_low_deg_inf} we immediately get the following corollary:
\begin{corollary}\label{corr:bound_low_deg_inf}
  There exists an absolute constant $C>0$ such that for all $\tau>0$,
    \[
    \tilde{W}_{k,d}[f]\leq C\cdot 2^{-k} d^{1-p}\tau^{2p-1} \Expect{x}{s_f(x)^p}
    +C\cdot \tau^{-\frac{1}{100 d}}\delta^{\frac{1}{200 d}}
    \left(\sum\limits_{\ell=1}^{\log d} 2^{-\card{\ell-k}}W_{\ell,d}[f]+\sum\limits_{\ell=1}^{\log n} 2^{-\card{\ell-k}}\eps_{\ell, d}\right).
    \]
\end{corollary}


  We may now give an upper bound on $\sum\limits_{k=1}^{\log d} W_{k,d}[f]$.
  \begin{claim}\label{claim:bd_outside_junta}
  For sufficiently large absolute constant $C_1$ and sufficiently large $C_2$ in comparison to $C_1$, we have that
    $\sum\limits_{k=1}^{\log d} W_{k,d}[f]\leq \frac{\eps}{16}$.
  \end{claim}
  \begin{proof}
    Let $\tau>0$ to be determined. As $W_{k,d}[f]\ll \tilde{W}_{k,d}[f]$, by Corollary~\ref{corr:bound_low_deg_inf} we have
    \[
    \sum\limits_{k=1}^{\log d} W_{k,d}[f]
    \ll
    \sum\limits_{k=1}^{\log d}
    \left(2^{-k}d^{1-p}\tau^{2p-1} A
    +\tau^{-\frac{1}{100 d}}\delta^{\frac{1}{200 d}}
    \left(\sum\limits_{\ell=1}^{\log d} 2^{-\card{\ell-k}}W_{\ell,d}[f]+\sum\limits_{\ell=1}^{\log n} 2^{-\card{\ell-k}}\eps_{\ell, d}\right)
    \right).
    \]
    The first sum is clearly $\ll d^{1-p}\tau^{2p-1} A$. The second sum is at most
    \[
    \tau^{-\frac{1}{100 d}}\delta^{\frac{1}{200 d}}
    \sum\limits_{\ell=1}^{\log d}W_{\ell,d}[f] \sum\limits_{k} 2^{-\card{\ell-k}}
    \ll \tau^{-\frac{1}{100 d}}\delta^{\frac{1}{200 d}}\sum\limits_{\ell=1}^{\log d} W_{\ell,d}[f].
    \]
    Finally, the third sum is at most
    \[
    \tau^{-\frac{1}{100 d}}\delta^{\frac{1}{200 d}}
    \sum\limits_{\ell=1}^{\log n} \eps_{\ell, d}
    \sum\limits_{k=1}^{\log d} 2^{-\card{\ell-k}}
    \ll
    \tau^{-\frac{1}{100 d}}\delta^{\frac{1}{200 d}}
    \sum\limits_{\ell=1}^{\log n} \eps_{\ell, d}
    \ll
    \tau^{-\frac{1}{100 d}}\delta^{\frac{1}{200 d}}
    W_{>d}[f]
    \ll \tau^{-\frac{1}{100 d}}\delta^{\frac{1}{200 d}}\eps,
    \]
    where we used~\eqref{eq:super_frid}.
    We thus get that
    \[
    \sum\limits_{k=1}^{\log d} W_{k,d}[f]
    \leq
    C\cdot d^{1-p}\tau^{2p-1} A
    +C\cdot \tau^{-\frac{1}{100 d}}\delta^{\frac{1}{200 d}}\sum\limits_{k=1}^{\log d} W_{k,d}[f]
    +C\cdot \tau^{-\frac{1}{100 d}}\delta^{\frac{1}{200 d}}\eps
    \]
    for some absolute constant $C>0$. We choose $\tau = \delta^{1/4}$ and then $C_2$ large enough so that
    $C\tau^{-\frac{1}{100 d}}\delta^{\frac{1}{200 d}}\leq \frac{1}{2}$, to get that
    $\sum\limits_{k=1}^{\log d} W_{k,d}[f]\leq 2C\cdot d^{1-p}\delta^{(2p-1)/4} A + 2C\delta^{\frac{1}{400 d}}\eps$,
    which is at most $\frac{\eps}{16}$ for large enough $C_2$.
  \end{proof}

  We can now prove Theorem~\ref{thm:super_fried}.
  \begin{proof}[Proof of Theorem~\ref{thm:super_fried}]
    Combining~\eqref{eq:super_frid1},~\eqref{eq:super_frid} and Claim~\ref{claim:bd_outside_junta} we get that
    $\sum\limits_{S: S\not\subseteq \mathcal{J}}\widehat{f}(S)^2\leq \frac{\eps}{8}$.
    Define $h(x) = 1$ if $\sum\limits_{S: S\subseteq \mathcal{J}}\widehat{f}(S)\chi_S(x)\geq 1/2$ and $0$ otherwise.
    Clearly, $h$ is a $\mathcal{J}$-junta and we have that
    \[
    \norm{f-h}_1
    =
    \norm{f-h}_2^2
    \leq 4\left\|f-\sum\limits_{S: S\subseteq \mathcal{J}}\widehat{f}(S)\chi_S(x)\right\|_2^2
    =4\sum\limits_{S: S\not\subseteq \mathcal{J}}\widehat{f}(S)^2\leq \eps.
    \]
    Finally, note that as the sum of $\left(1-\frac{1}{3d}\right)$-noisy influences of $f$ is at most $O(d)$ by Fact~\ref{fact:sum_of_noisy}, we get that
    \[
    \card{\mathcal{J}}\leq O\left(\frac{d}{\delta}\right).
    \qedhere
    \]
  \end{proof}

\subsection{Proof of Claim~\ref{claim:super_fried_main}}\label{sec:missing_proofs}
For the proof of Claim~\ref{claim:super_fried_main} we need a version of the hypercontractive inequality with the noise operator, as follows:
\begin{lemma}\label{lemma:hypercontractivity_noise}[\cite{ODonnell}]
  Let $q>2$, and $\rho\leq \frac{1}{\sqrt{q-1}}$. Then for all $f\colon\{0,1\}^n\to\mathbb{R}$ we have
  $\norm{T_{\rho} f}_q\leq \norm{f}_2$.
\end{lemma}

We also need the following claim, asserting that noise only decreases Talagrand's boundary.
\begin{claim}\label{claim:noise_decreases_tal}
  For $f\colon \{0,1\}^n\to\mathbb{R}$ and any noise operator $S = T_{\rho_1}\otimes\ldots T_{\rho_n}$ we have
  $\Expect{x}{\norm{\nabla{S f}(x)}_p}\leq \Expect{x}{\norm{\nabla{f}(x)}_p}$.
\end{claim}
\begin{proof}
  \[\Expect{x}{\norm{\nabla (S f)(x)}_p}
  \leq
  \Expect{x}{\norm{S \nabla f(x)}_p}
  =\Expect{x}{\norm{\Expect{y\sim S x}{\nabla f(y)}}_p},
  \]
  and the result follows from the triangle inequality.
\end{proof}


  Take $(J,z)$ to be a random restriction so that each $j\in [n]$ is included in $J$ independently with probability $\frac{1}{3d}$,
  and consider $g = S_k'' f$.
  We calculate
  the contribution of $\bar{\mathcal{J}}$ to the level $1$ weight of a restriction of $g$:
  \begin{align*}
  d\Expect{J, z}{\sum\limits_{j\in J\cap \bar{\mathcal{J}}}\widehat{g_{\overline{J}\rightarrow z}}(\set{j})^2}
  =d\sum\limits_{j\in \overline{\mathcal{J}}}\Expect{J}{\sum\limits_{S}\widehat{g(S)}^2 1_{S\cap J = \{j\}}}
  &=d\sum\limits_{S}\frac{\card{S\cap \bar{\mathcal{J}}}}{3d}\left(1-\frac{1}{3d}\right)^{\card{S}-1}\widehat{g(S)}^2\\
  &=\frac{1}{3}\sum\limits_{S}\card{S\cap \bar{\mathcal{J}}}\left(1-\frac{1}{3d}\right)^{\card{S}-1}\widehat{g(S)}^2.
  \end{align*}
  Also,
  \begin{align*}
    \sum\limits_{j\not\in \mathcal{J}} \norm{S_k \partial_j f}_2^2
    =
    4\sum\limits_{S}
    \card{S\cap \bar{\mathcal{J}}}
    \left(1-\frac{1}{2^k}\right)^{-2}\rho_1^{2\card{S\cap \mathcal{J}}}
    \rho_2^{2\card{S\cap \bar{\mathcal{J}}}}\widehat{g(S)}^2,
  \end{align*}
  where $\rho_1 = \frac{1-1/d}{1-2/(3d)}$ and $\rho_2 = \frac{1-1/2^k}{1-1/2^k+1/(3d)}$. Using $\frac{a}{b}\leq \frac{a+c}{b+c}$ for $0\leq a\leq b$ and $c\geq 0$ we have
  \[
  \rho_1 \leq \frac{1-1/d+2/(3d)}{1-2/(3d)+2/(3d)} = 1-\frac{1}{3d},
  \qquad
  \rho_2 \leq \frac{1-1/2^k+(1/2^k-1/(3d))}{1-1/2^k+1/(3d)+(1/2^k-1/(3d))} = 1-\frac{1}{3d},
  \]
  and so
  \[
  \sum\limits_{j\not\in \mathcal{J}} \norm{S_k \partial_j f}_2^2\leq
  12\left(1-\frac{1}{d}\right)^{-2} d\Expect{J, z}{\sum\limits_{j\in J\cap \bar{\mathcal{J}}}\widehat{g_{\overline{J}\rightarrow z}}(\set{j})^2}
  \ll d\Expect{J, z}{\sum\limits_{j\in J\cap \bar{\mathcal{J}}}\widehat{g_{\overline{J}\rightarrow z}}(\set{j})^2}.
  \]
  Hence our
  goal of upper bounding the left hand side in Claim~\ref{claim:super_fried_main} reduces to
  upper bounding the contribution of $\mathcal{J}$ to the level $1$ of random restrictions of $g$. For $J, z$ define
 \[
  L_{J,z} = \sett{j\in J\cap \bar{\mathcal{J}}}{\card{\widehat{g_{\overline{J}\rightarrow z}}(\set{j})} \leq \tau},
  \qquad
  H_{J,z} = \sett{j\in J\cap \bar{\mathcal{J}}}{\card{\widehat{g_{\overline{J}\rightarrow z}}(\set{j})} > \tau}.
  \]
  Then
  \[
  \Expect{J, z}{\sum\limits_{j\in J\cap \bar{\mathcal{J}}}\widehat{g_{\overline{J}\rightarrow z}}(\set{j})^2}
  =\Expect{J, z}{\sum\limits_{j\in L_{J,z}}\widehat{g_{\overline{J}\rightarrow z}}(\set{j})^2}
  +\Expect{J, z}{\sum\limits_{j\in H_{J,z}}\widehat{g_{\overline{J}\rightarrow z}}(\set{j})^2}.
  \]

  \paragraph{Bounding the contribution from $L_{J,z}$.}
  For the first sum, we have
  \begin{align*}
  \sum\limits_{j\in L_{J,z}}\widehat{g_{\overline{J}\rightarrow z}}(\set{j})^2
  \leq
  \left(\sum\limits_{j\in L_{J,z}}\widehat{g_{\overline{J}\rightarrow z}}(\set{j})^2\right)^{p}
  &\leq \left(\sum\limits_{j\in L_{J,z}}\tau^{(2p-1)/p}\card{\widehat{g_{\overline{J}\rightarrow z}}(\set{j})}^{1/p}\right)^{p}\\
  &\leq \tau^{2p-1}\left(\sum\limits_{j}\card{\widehat{g_{\overline{J}\rightarrow z}}(\set{j})}^{1/p}\right)^{p}.
  \end{align*}
  In the last expression we have $\tau^{2p-1}$ times the $1/p$-norm of the vector $(\widehat{g_{\overline{J}\rightarrow z}}(\set{j}))_{j\in J}$.
  Thus, using the triangle inequality in a similar way to Lemma~\ref{lemma:tal_lvl1}, we may bound the last expression from above by
  \begin{equation}\label{eq3}
  \tau^{2p-1}\Expect{x}{\norm{\nabla{g_{\overline{J}\rightarrow z}}(x)}_{1/p}}.
  \end{equation}
  Taking expectation, we get that
  \[
  \Expect{J,z}{\sum\limits_{j\in L_{J,z}}\widehat{g_{\overline{J}\rightarrow z}}(\set{j})^2}
  \leq \tau^{2p-1}\Expect{J,z}{\Expect{x}{\norm{\nabla {g_{\overline{J}\rightarrow z}}(x)}_{1/p}}}
  = \tau^{2p-1}\Expect{x}{\Expect{J,z}{\norm{\nabla g_{\overline{J}\rightarrow z}(x)}_{1/p}}}.
  \]
  Fix $x,z$ and consider the vector $\nabla g(x,z)$.
  Note that
  \[
  \Expect{J}{\norm{\nabla g_{\overline{J}\rightarrow z}(x)}_{1/p}^{1/p}}
  =
  \Expect{J}{\sum\limits_{j}\card{\partial_j g(x,z)}^{1/p}1_{j\in J}}
  =\frac{1}{3d}\norm{\nabla g(x,z)}_{1/p}^{1/p},
  \]
  and it follows by Jensen's inequality that
  \[
  \Expect{J}{\norm{\nabla g_{\overline{J}\rightarrow z}(x)}_{1/p}}
  \leq
  \Expect{J}{\norm{\nabla g_{\overline{J}\rightarrow z}(x)}_{1/p}^{1/p}}^{p}
  \leq d^{-p} \norm{\nabla g(x,z)}_{1/p}.
  \]
  It follows that~\eqref{eq3} is upper bounded by $\tau^{2p-1}d^{-p}\Expect{x,z}{\norm{\nabla g(x,z)}_{1/p}}$.
  Using Claim~\ref{claim:noise_decreases_tal}, we have that this may be upper bounded by
  $\tau^{2p-1}d^{-p}\Expect{x,z}{\norm{\nabla f(x,z)}_{1/p}} = \tau^{2p-1}d^{-p} \Expect{x,z}{s_f(x,z)^p}$.



  \paragraph{Bounding the contribution from $H_{I,z}$.}
  We have
  \[
  \sum\limits_{j\in H_{J,z}}\widehat{g_{\overline{J}\rightarrow z}}(\set{j})^2
  \leq \tau^{-\frac{1}{100d}}\sum\limits_{j\in H_{J,z}}\card{\widehat{g_{\overline{J}\rightarrow z}}(\set{j})}^{2+\frac{1}{100d}}
  \leq \tau^{-\frac{1}{100d}}\sum\limits_{j\in J\cap\bar{\mathcal{J}}}\card{\widehat{g_{\overline{J}\rightarrow z}}(\set{j})}^{2+\frac{1}{100d}},
  \]
  and we next take expectation over $z$:
  \begin{equation}\label{eq5}
  \Expect{z}{\sum\limits_{j\in H_{J,z}}\widehat{g_{\overline{J}\rightarrow z}}(\set{j})^2}
  \leq
  \tau^{-\frac{1}{100d}}\sum\limits_{j\in J\cap\bar{\mathcal{J}}}\Expect{z}{\card{\widehat{g_{\overline{J}\rightarrow z}}(\set{j})}^{2+\frac{1}{100d}}}.
  \end{equation}
  Define the operator
  $T = T_{\rho_1}^{\mathcal{J}}T_{\rho_2}^{\overline{\mathcal{J}}}$
  where $\rho_1 = \frac{1-2/(3d)}{1-1/(3d)}$ and $\rho_2 = \frac{1-1/2^k+1/(3d)}{1-1/2^k+2/(3d)}$,
  and note that by the semi-group property of the noise operator (namely, the fact that $T_{\rho}T_{\rho'} = T_{\rho\rho'}$)
  we may write $g = T S_k' f$. Thus, $\widehat{g_{\overline{J}\rightarrow z}}(\set{j})$ is equal to
  \[
  \sum\limits_{S\subseteq \overline{J} \cup\{j\}, S\ni j} \widehat{g}(S)\chi_{S\setminus\{j\}}(z) =
  \sum\limits_{S\subseteq \overline{J} \cup\{j\}, S\ni j}
  \rho_1^{\card{S\cap \mathcal{J}}}
  \rho_2^{\card{S\cap \overline{\mathcal{J}}}}
  \widehat{S_k' f}(S)\chi_{S\setminus\{j\}}(z).
  \]
  Thus, using hypercontractivity, i.e.~Lemma~\ref{lemma:hypercontractivity_noise}, we have that
  \begin{align}\label{eq4}
  \Expect{z}{\card{\widehat{g_{\overline{J}\rightarrow z}}(\set{j})}^{2+\frac{1}{100d}}}
  &=\left\|\sum\limits_{S\subseteq \overline{J} \cup\{j\}, S\ni j}
  \rho_1^{\card{S\cap \mathcal{J}}}
  \rho_2^{\card{S\cap \overline{\mathcal{J}}}}
  \widehat{S_k' f}(S)\chi_{S\setminus\set{j}}(z)\right\|_{2+\frac{1}{100d}}^{2+\frac{1}{100d}}\notag\\
  &\leq \left\|\sum\limits_{S\subseteq \overline{J} \cup\{j\}, S\ni j}
  \rho_1^{\card{S\cap \mathcal{J}}}
  \rho_2^{\card{S\cap \overline{\mathcal{J}}}}
  \left(1-\frac{1}{3d}\right)^{-(\card{S}-1)}\widehat{S_k' f}(S)\chi_{S\setminus\set{j}}(z)\right\|_2^{2+\frac{1}{100d}}\notag\\
  &\leq \left\|\sum\limits_{S\ni j}\widehat{S_k' f}(S)\chi_{S\setminus\set{j}}(z)\right\|_2^{2+\frac{1}{100d}}.
  \end{align}
  In the last inequality, we used Parseval's equality and that fact that $\rho_1,\rho_2\leq 1-\frac{1}{3d}$.
  %
  %
%
%
%
%
%
%
%
%
%
%
  Note that
  \[
  \left\|\sum\limits_{S\ni j}\widehat{S_k' f}(S)\chi_{S\setminus\set{j}}(z)\right\|_2^{2+\frac{1}{100d}}
  \ll \norm{S_k' \partial_j f}_2^{2+\frac{1}{100d}},
  \]
  so plugging~\eqref{eq4} into~\eqref{eq5} yields that
  \[
  \Expect{z}{\sum\limits_{j\in H_{J,z}}\widehat{g_{\overline{J}\rightarrow z}}(\set{j})^2}\ll
  \tau^{-\frac{1}{100d}}\sum\limits_{j\in J\cap\bar{\mathcal{J}}}\norm{S_k' \partial_j f}_2^{2+\frac{1}{100d}}.
  \]
  Taking expectation over $J$ now gives that
  \[
  \Expect{z,J}{\sum\limits_{j\in H_{J,z}}\widehat{g_{\overline{J}\rightarrow z}}(\set{j})^2}
  \ll
  \tau^{-\frac{1}{100d}}\sum\limits_{j\in\bar{\mathcal{J}}} \norm{S_k' \partial_j f}_2^{2+\frac{1}{100d}} \Expect{J}{1_{j\in J}}
  = \frac{1}{3d}\tau^{-\frac{1}{100d}}\sum\limits_{j\in\bar{\mathcal{J}}} \norm{S_k' \partial_j f}_2^{2+\frac{1}{100d}}.
  \]
\qed
